\documentclass[runningheads]{llncs}

% ---------------------------------------------------------------
% Include basic ECCV package
\usepackage{eccv}

% ---------------------------------------------------------------
% Other packages
\usepackage{eccvabbrv}
\usepackage{graphicx}
\usepackage{booktabs}
\usepackage{multirow}
\usepackage{algorithm}
\usepackage{algpseudocode}

% Keep absolute values aligned and append a compact relative delta on the right.
\newcommand{\scoredelta}[2]{\makebox[2.8em][r]{#1}\hspace{0.06em}{\scriptsize(#2)}}

% Improves PDF readability for those with disabilities.
\usepackage[accsupp]{axessibility}

% ---------------------------------------------------------------
% Hyperref package
\usepackage{hyperref}

% ---------------------------------------------------------------
% Counter Reset and Prefix Configuration
\setcounter{section}{0}
\setcounter{figure}{0}
\setcounter{table}{0}
\setcounter{equation}{0}
\setcounter{algorithm}{0}

\renewcommand{\thesection}{S\arabic{section}}
\renewcommand{\thefigure}{S\arabic{figure}}
\renewcommand{\thetable}{S\arabic{table}}
\renewcommand{\theequation}{S\arabic{equation}}
\renewcommand{\thealgorithm}{S\arabic{algorithm}}

\begin{document}

% ---------------------------------------------------------------
\title{OREO: Fidelity Alignment in 3D Generation via On-the-fly Rendering-Editing Optimization\\---Supplementary Material---}

\titlerunning{OREO: Supplementary Material}

\author{}
\institute{}

\maketitle

\section{Overview}
In this supplementary material, we provide:
\begin{itemize}
  \item Section~\ref{sec:implementation_details}: Dataset construction, training details, and editing setup for compared methods.
  \item Section~\ref{sec:user_study_details}: Comprehensive protocols and setup for our user preference study.
  \item Section~\ref{sec:computational_overhead}: Computational overhead of the editing component during OREO post-training.
  \item Section~\ref{sec:limitations_failure}: Discussion on failure cases and limitations.
\end{itemize}

% ---------------------------------------------------------------
\section{Implementation and Dataset Details}
\label{sec:implementation_details}
This section provides the implementation details used in our main experiments, including dataset construction, training hyper-parameters, and the editing setup for the compared methods.

\subsection{Conceptual Design Dataset Construction}
\label{sec:dataset_construction}
We construct an image-only collection to support generator post-training without paired 3D supervision.
The dataset comprises 2,396 concept-design reference images that we manually curate from public online sources; most images are AI-generated or user-shared creations posted on public content platforms, reflecting the imaginative and stylized inputs that image-to-3D users bring in practice.
We prioritize references with clean foregrounds, well-defined silhouettes, and informative texture or material cues, and manually filter the collection to remove duplicates, low-resolution thumbnails, watermarked images, and content with unclear foregrounds.
We will release the dataset on our project page.

\subsection{Training Details}
\label{sec:training_details}
We post-train the Trellis-image-large 3D generator with OREO on the Conceptual Design Dataset training split, using Qwen-Image-Edit-2511~\cite{wu2025qwen} as the frozen 2D editor.

\textbf{Optimization.}
We use the Adan optimizer with learning rate $10^{-4}$, weight decay $0$, and $\epsilon=10^{-4}$.
Training runs in \texttt{bf16} mixed precision, with a per-GPU batch size of $1$ and no gradient accumulation, giving an effective batch size of $8$ reference images per iteration.
We train for $10$ epochs over the Conceptual Design Dataset training split.
The full-parameter update mode is applied to the sparse-latent flow-matching branch of Trellis, while the sparse-structure branch and the decoder $\mathcal{D}$ remain frozen.

\textbf{Rollout and camera sampling.}
Each iteration performs an on-policy ODE rollout with the current student generator to obtain a clean latent $z_0 = G_\theta(x^{\texttt{ref}})$.
For every reference we sample a single training view with yaw uniformly drawn from $[0^\circ, 360^\circ]$, pitch fixed at $0^\circ$, camera distance $r=2.0$, and field of view $40^\circ$.
The camera distance is further adjusted by an adaptive-distance rule (fill ratio $0.9$) so that the object occupies a consistent silhouette across scales.
Rendering uses the Gaussian-splatting renderer at $1024\times1024$ resolution with a white background.

\textbf{Loss and time-step sampling.}
For the contrastive objective in Eq.~13 of the main paper, the training time step is sampled from $t \sim \mathcal{U}(0.02, 0.98)$, and the pretrained-generator forward passes for $z_0^{+}$ and $z_0^{-}$ share the same fresh noise draw $\epsilon'$ and are computed without classifier-free guidance, i.e., a single conditional forward pass under $x^{\texttt{tgt}}$ and $x^{\texttt{src}}$ respectively.
The total loss is $\mathcal{L}_{total} = \mathcal{L}_{contrast} + \lambda\,\mathcal{L}_{reg}$, with $\lambda = 1.0$.
The velocity regularizer $\mathcal{L}_{reg}$ compares the student and pretrained velocity fields at the same $(z_t, t, x^{\texttt{ref}})$.

\subsection{Editing Setup for Compared Methods}
\label{sec:editing_setup}
Reinforced Editing and the off-the-shelf 2D editors used as baselines share the same underlying editor backbone but differ in their text-conditioning interface.
Both prompt families are fixed across the entire dataset to avoid per-example prompt tuning.

\textbf{Reinforced Editing.}
We instantiate Reinforced Editing on top of Qwen-Image-Edit-2511 with $N=12$ total time steps and $N_e=9$ editing steps applied at the last portion of the schedule (editing ratio $0.75$).
The target branch uses classifier-free guidance with scale $w=4$, while the source branch is an unconditional forward pass, in accordance with Eq.~7 of the main paper.
The shared noise $\epsilon$ is initialized once per view and updated at each step by injecting the CFG signal, as in Eq.~8 of the main paper.
The editing resolution matches the renderer at $1024\times1024$.
The text component of $c^{\texttt{ref}}$ is:
\begin{quote}
\itshape
``Rotate the camera. White background.''
\end{quote}

\textit{Prompt selection.}
We arrived at this two-phrase instruction after sweeping alternative prompts on the Conceptual Design Dataset, holding the editor and all other RE hyper-parameters fixed.
Table~\ref{tab:prompt_ablation} groups representative candidates by intent.
Fidelity-only prompts do not disambiguate the viewpoint role of $x^{\texttt{src}}$ and yield only modest CLIP/DINO gains; underspecified viewpoint prompts (e.g., \textit{Move the camera}) improve on this but still trail the ``rotate'' formulation; asking for a full 3D reinterpretation (\textit{Generate a 3D model of the image}) actively degrades fidelity, turning the CLIP delta negative.
``\textit{Rotate the camera}'' strikes the best balance because it signals a novel view of the \emph{same} object without over-specifying azimuth or elevation, which we found in preliminary trials to trigger hallucinated backgrounds around the requested viewpoint.
Adding ``\textit{White background}'' further improves both deltas by aligning the pseudo-GT with the 3D backbone's white-background rendering convention; without it, the editor occasionally fills in scene backgrounds absent from $x^{\texttt{src}}$ (studio lighting, ground planes, vegetation), leaking pixel-level supervision into the generator.
Additional stylistic modifiers on top (\textit{Consistent Lightening}, \textit{visual style}, \textit{concept design}) do not consistently help, so we adopt the minimal two-phrase form.

\begin{table}[t]
  \centering
  \caption{Prompt selection for Reinforced Editing on the Conceptual Design Dataset. All entries share the same editor and RE hyper-parameters ($N=12$, $N_e=9$, $w=4$); we report $\Delta$CLIP / $\Delta$DINO similarity to $x^{\texttt{ref}}$ relative to the unedited $x^{\texttt{src}}$, and Mask IoU with $x^{\texttt{src}}$. \textbf{Bold} row is our final choice.}
  \label{tab:prompt_ablation}
  \setlength{\tabcolsep}{4pt}
  \begin{tabular}{@{}lccc@{}}
    \toprule
    Prompt & $\Delta$CLIP $\uparrow$ & $\Delta$DINO $\uparrow$ & IoU $\uparrow$ \\
    \midrule
    \multicolumn{4}{@{}l}{\textit{Fidelity-only}} \\
    Super resolution of the input image           & +0.0209 & +0.0371 & 0.9398 \\
    Enrich the detail given the reference image   & +0.0171 & +0.0250 & 0.9469 \\
    Add the appearance of image 1                 & +0.0233 & +0.0346 & 0.9406 \\
    \midrule
    \multicolumn{4}{@{}l}{\textit{Viewpoint change}} \\
    Move the camera                                & +0.0277 & +0.0398 & 0.9406 \\
    Obtain another side view                       & +0.0250 & +0.0341 & 0.9330 \\
    Generate a novel view                          & +0.0161 & +0.0264 & 0.9437 \\
    Move the camera to a novel view                & +0.0268 & +0.0387 & 0.9408 \\
    \midrule
    \multicolumn{4}{@{}l}{\textit{Full 3D reinterpretation}} \\
    Generate a 3D model of the image               & $-0.0274$ & +0.0062 & 0.9455 \\
    \midrule
    \multicolumn{4}{@{}l}{\textit{``Rotate the camera'' variants}} \\
    Rotate the camera                              & +0.0314 & +0.0371 & 0.9402 \\
    Rotate the camera. Consistent Lightening       & +0.0204 & +0.0357 & 0.9417 \\
    Rotate the camera. Consistent visual style     & +0.0206 & +0.0310 & 0.9428 \\
    Rotate the camera. Consistent concept design   & +0.0265 & +0.0316 & 0.9422 \\
    \textbf{Rotate the camera. White background.}\ (Ours) & \textbf{+0.0339} & \textbf{+0.0387} & \textbf{0.9411} \\
    \bottomrule
  \end{tabular}
\end{table}

\textbf{Off-the-shelf editor baselines.}
For the feedback comparison in Sec.~4.1 of the main paper, we compare against the original Qwen-Image-Edit-2511 pipeline and NanoBanana Pro.
Because both editors expose a single visual context that concatenates the source view $x^{\texttt{src}}$ and the reference $x^{\texttt{ref}}$ side by side, the prompt must first assign roles to the two images before requesting the edit:
\begin{quote}
\itshape
``The first image is a 3D rendered view. The second image is the reference. Edit the first image to match the reference object's appearance. White background.''
\end{quote}
The prompt is a minimal adaptation of the effective instruction reported by Photo3D and is applied identically to both editors.
Only the underlying editor backbone is swapped between runs; API-specific parameters such as sampler, classifier-free guidance scale, and number of denoising steps are listed in the released source code.

% ---------------------------------------------------------------
\section{User Preference Study Protocols}
\label{sec:user_study_details}
We conduct a user preference study to complement the automatic metrics in the main paper.
The study focuses on whether the generated 3D assets preserve the identity of the input image while improving visual fidelity across multiple rendered views.

\textbf{Study setup.}
We use the same 100 held-out examples from the Conceptual Design Dataset as in the main evaluation.
For each example, participants are shown the input reference image and multi-view renderings generated by three methods: Pretrained (Trellis), Photo3D, and OREO.
The method names are hidden during evaluation, and participants are asked to choose the result with the best overall quality.

\textbf{Evaluation criteria.}
Participants are instructed to consider three aspects jointly.
First, \textit{Fidelity} measures whether the generated asset contains realistic materials, sharp boundaries, and fine-grained texture details.
Second, \textit{Consistency} measures whether the rendered views remain coherent as a 3D object rather than showing view-dependent artifacts.
Third, \textit{Identity} measures whether the generated asset preserves the semantic identity and distinctive details of the input reference image.

\textbf{Participants.}
The study includes 20 participants, including participants familiar with 3D reconstruction and computer graphics.
All participants evaluate the same 100 examples, and we aggregate their preferences over all responses.

\begin{table}[t]
  \centering
  \caption{User preference study on the held-out 100-example Conceptual Design Dataset. Participants compare multi-view renderings with respect to reference fidelity, 3D consistency, and identity preservation.}
  \label{tab:user_preference}
  \begin{tabular}{@{}lc@{}}
    \toprule
    Method & Preference Ratio $\uparrow$ \\
    \midrule
    Pretrained (Trellis) & 29\% \\
    Photo3D & 33\% \\
    OREO (Ours) & \textbf{38\%} \\
    \bottomrule
  \end{tabular}
\end{table}

As shown in Table~\ref{tab:user_preference}, OREO obtains the highest preference ratio.
This indicates that human evaluators more frequently prefer OREO's balance between reference fidelity, fine-grained visual details, and multi-view 3D consistency.
Since only aggregate preference ratios are recorded, we report the descriptive preference results and do not include unsupported significance tests.

% ---------------------------------------------------------------
\section{Computational Overhead}
\label{sec:computational_overhead}
Each OREO training iteration consists of two costs: a generator update (rollout, render, and backward pass), and a call to the frozen 2D editor to produce the pseudo-GT $x^{\texttt{tgt}}$.
The generator update alone takes roughly 22~s per iteration.
If the 9-step editing call is executed sequentially inside the same iteration, it dominates the wall-clock cost, contributing about 81\% of the total per-iteration time and stretching one iteration to roughly 5$\times$ the generator-only cost.
To avoid paying this full cost, we wrap the editor as an asynchronous service that runs in parallel with the generator update on the next batch, which reduces the residual editor time exposed to the training loop to about 45\% of a training iteration, with no observed change in optimization dynamics or final quality.
Memory usage is dominated by the generator, and we optionally enable activation checkpointing on the 3D generator when running on memory-constrained hardware.

\begin{table}[t]
  \centering
  \caption{Per-iteration cost breakdown for OREO post-training. \emph{Effective editor wait} is the residual editor time not overlapped by the concurrent generator update.}
  \label{tab:computational_overhead}
  \begin{tabular}{@{}lcp{0.42\linewidth}@{}}
    \toprule
    Component & Cost & Note \\
    \midrule
    Generator update & $\sim$22~s/iter & Rollout, render, forward/backward. \\
    Editor call (sequential) & $\sim$81\% of iter time & Dominates iteration if run in-line. \\
    Editor call (async, effective wait) & $\sim$45\% of iter time & Overlapped with next-batch generator update. \\
    \bottomrule
  \end{tabular}
\end{table}

% ---------------------------------------------------------------
\section{Limitations and Failure Cases}
\label{sec:limitations_failure}
OREO inherits its appearance guidance from the 2D editing prior, and thus may also inherit the editor's compositional or viewpoint biases in rare cases.
We identify three representative failure modes on the Conceptual Design Dataset and describe each below with a short mechanistic note.

\textbf{Viewpoint drift on face-like objects.}
For face-like objects or characters, the 2D editor may rotate the face toward the camera instead of preserving the intended 3D-facing direction, leading to view-dependent orientation inconsistency.
This drift arises from a mismatch between the 2D editor's canonical portrait prior and the 3D view-consistency requirement, so pseudo-GTs may reorient a face toward the camera; repeated use of such pseudo-GTs can improve local facial details while causing inconsistent face orientation across views.

\textbf{Editor-inherited color and style bias.}
The dynamic noise update tightens alignment with the editor's appearance prior and can therefore inherit editor-specific color bias on unusual materials such as stone or metallic surfaces, shifting the intended hue while preserving geometry.
Similar biases can also appear as style shifts on stylized references whose material is far from the editor's photographic training distribution.

\textbf{Remaining artifacts on weakly constrained regions.}
Rarely observed regions such as occluded surfaces, thin appendages, and sharp material boundaries receive weaker per-view supervision and can retain small local artifacts on back views and along thin structures.
These regions are also under-constrained by the reference $x^{\texttt{ref}}$, which only depicts a single canonical viewpoint.

\textbf{Reliance on jointly-modeled geometry and appearance.}
OREO is designed for 3D generators such as Trellis whose latent space jointly encodes geometry and appearance, so a single latent update can simultaneously refine texture detail while keeping structure stable through the velocity regularizer.
For pipelines that decouple geometry and texture, e.g., generating a mesh first and then baking a texture map in a separate stage, the contrastive signal from an edited RGB view can only supervise the texture stage and is largely disconnected from the geometry decoder.
In this decoupled setting, we expect the fidelity gain to be smaller and, more importantly, structural consistency between the geometry and the injected appearance is no longer regularized by the same optimizer.
Extending OREO to decoupled pipelines therefore requires either a texture-only variant of the contrastive objective or an additional geometry-aware regularizer on the mesh stage; we leave this as future work.

These observations suggest future directions such as stronger multi-view consistency constraints, 3D-aware editing models, or uncertainty-aware filtering of pseudo-GTs before they are used for generator updates.

% ---------------------------------------------------------------
% References
\bibliographystyle{splncs04}
\bibliography{main}

\end{document}
